\section{The predictable acquired-geometry theorem}
\label{sec:main}
For a dictionary $C=(r_1,\ldots,r_m)\in H^m$ put
\[
 d_C(x)=w(x)+\min_j\|f(x)-r_j\|^2.
\]
The acquired marked profile is
\begin{equation}\label{eq:profile}
 \Qm_m(b)=\min\left\{\int d_C h\,d\nu:
       C\in H^m,\ 0\le h\le b,\ \int h\,d\nu=1\right\},\qquad b\ge1.
\end{equation}
It is $+\infty$ for $b<1$. Here $h\nu$ is a candidate selected \emph{probability measure}, not a free preparation. Its actual acquisition cost enters through $b$. The minimum is over Borel representatives modulo $\nu$-null sets.

\begin{theorem}[Predictable acquired geometry with paid causal acquisition]\label{thm:main}
Consider the raw experiment of Section~\ref{sec:setup}. Assume the actual-fiber criterion \eqref{eq:fiber-ratio}, or an independently verified comparison \eqref{eq:task-compare}; the actual conditional mass bounds \eqref{eq:mass-upper}--\eqref{eq:mass-lower}; the lower effort certificate \eqref{eq:effort-drift}; and the $B$-live-label program $\mathcal G$. Set
\[
 b_-(N)=uN/\kappa_-,\qquad
 b_+(N)=\max\{1,\ell N/\kappa_+\}.
\]
For $N\ge\kappa_+$ and $m\ge1$,
\begin{equation}\label{eq:cp-main}
 k_-\Qm_m(b_-(N))\le R_\cp^\Theta(N,m)-B_0
       \le k_+\frac u\ell\Qm_m(b_+(N)).
\end{equation}
For $M\ge B+1$,
\begin{equation}\label{eq:aut-main}
 k_-\Qm_{M-1}(b_-(N))\le R_\aut^\Theta(N,M)-B_0
       \le k_+\frac u\ell\Qm_{M-B}(b_+(N)).
\end{equation}
Every finite upper is implemented at the original budget by one common Borel routing program with exactly the stated upper alphabet bound. Proposition~\ref{prop:finite-cells} supplies a finite-description implementation under its explicit computational certificate. No exact finite observable primary driver is required. No interchange of worst-model and best-program operations is used.

If every legal autonomous procedure requires at least $b_0$ live labels under one common admitted model, the lower profile in \eqref{eq:aut-main} is $\Qm_{M-b_0}$. In particular, if $b_0=B$, the same paid terminal width occurs on both sides. Section~\ref{sec:memory} supplies a physical-command certificate for this condition. No procedure is feasible for $N<\kappa_-$. All profile minima at finite $b\ge1$ exist.
\end{theorem}

The two resource arguments are intentional. Bounded conditional density ratios do not imply constant-factor comparison of $\Qm(cb)$ and $\Qm(b)$ at a critical atom. Theorem~\ref{thm:atomic} preserves that transition. Under regular acquired geometry, Corollary~\ref{cor:regular} gives a matching law. Thus the theorem supplies both an all-strategy lower and a finite-state upper, rather than identifying a relaxation with an implementation by definition.

A physical experiment may contain several additive resources. The same construction bounds the expected amount of resource $j$ by $\kappa_{+,j}b/\ell$ whenever the acquisition program has that conditional attempt-cost bound. To obtain an upper satisfying all budgets, choose $b$ using their minimum admissible value. A lower can use any nonnegative linear combination of resources with a verified drift inequality. One constraint with a positive per-call lower bound supplies finite expected duration; bytes, raw calls and physical time are not summed without explicit conversion weights.

\begin{corollary}[Equivalence of certified task marks]\label{cor:mark-invariance}
If two bounded marks $(f,w)$ and $(\widetilde f,\widetilde w)$ on the same observed mark space satisfy uniform comparisons with the same full conditional task, their marked distortions are comparable for every prediction, and their acquired profiles are comparable with the same $b,m$. In particular any such mark can replace \eqref{eq:canonical-mark}, with explicit changes only in the comparison constants.
\end{corollary}
\begin{proof}
For each observed $x$ choose any attainable conditional pair on its nonempty task fiber and compare both marked distortions to its risk. The inequalities hold for every $r$, hence also for the minimum over a common finite dictionary. Integration against any common $h\nu$ and then minimization prove the profile comparison. No change of preparation or renormalization of a stratum is involved.
\end{proof}
